\section{The Shape of the Taxonomy}

\subsection{Two distinctions, not four boxes}

The familiar labels tempt the mind to a two-by-two grid: divine or human crossed with natural or positive, yielding four equal squares. The classical doctrine does not draw that grid, and seeing why is most of the payoff of the preceding three sections.

The \emph{divine/human} distinction divides by \textbf{author}: whose reason and will does this law proceed from? The \emph{natural/positive} distinction divides by \textbf{mode of institution and promulgation}: is this law given in the very constitution of the rational creature, knowable by the light of reason (promulgated, as q.~90 a.~4 ad 1 says, by God's instilling it into human minds), or is it enacted by an act in history---spoken, written, promulgated to a community at a time?

Cross them honestly and the grid is asymmetric. \emph{Divine natural} law exists: it is the natural law itself, whose author is God legislating through creation. \emph{Divine positive} law exists: the Old and New Law, revealed in history. \emph{Human positive} law exists: civil and ecclesiastical enactment and custom. But the fourth square---\emph{human natural} law---is empty, and necessarily so: no human authority authors nature. What occupies the place people reach for when they say such things is human \emph{recognition, formulation, and determination} of natural law---the conclusions and determinations of q.~95 a.~2---which are acts about the natural law, not a second natural law. And standing over the whole scheme, not inside it, is the eternal law: not a fifth species alongside the others but the divine wisdom of which natural law is the participation in rational creatures (q.~91 a.~2) and from which every true law, including human law, derives whatever binding force it has (q.~93 a.~3; q.~96 a.~4).\endnote{This paragraph is the article's synthesis of the loci quoted in sections 3--5; the asymmetry claim itself---no fourth, humanly authored ``natural'' law---is standard school doctrine rather than a novelty, but its formulation here is the article's own and should be read as interpretation of the quoted texts, not as a further quotation.}

One more asymmetry matters juridically. Within human positive law, the Church's tradition distinguishes law by \emph{which community's} care it serves: civil law, enacted by the civil community for temporal common good; and ecclesiastical law, enacted by the Church's authority for the common good of the faithful ordered to salvation. The 1983 Code presupposes both, defers to civil law in defined matters (cc.~22, 1290), and calls its own purely human enactments \latin{leges mere ecclesiasticae}---\emph{merely} ecclesiastical laws (cc.~11, 85)---precisely to mark them off from the divine law, natural and positive, that the same Code also carries but did not author.

\subsection{The taxonomy in one table}

The table below states the scheme as the remainder of this article will use it. It is a summary of the sources already quoted, with the canonical consequences that sections 8--10 will verify canon by canon; the ``exemplary loci'' column names where each row's claims are documented in this article.

\begingroup
\footnotesize
\begin{longtable}{>{\bfseries\raggedright\arraybackslash}p{0.115\linewidth}>{\raggedright\arraybackslash}p{0.155\linewidth}>{\raggedright\arraybackslash}p{0.17\linewidth}>{\raggedright\arraybackslash}p{0.15\linewidth}>{\raggedright\arraybackslash}p{0.17\linewidth}>{\raggedright\arraybackslash}p{0.145\linewidth}}
\toprule
\textbf{Kind of law} & \textbf{Author / source} & \textbf{How instituted and promulgated} & \textbf{Mutability} & \textbf{Dispensability / custom} & \textbf{Exemplary loci}\\
\midrule
\endhead
Eternal law & God; the plan of divine wisdom directing all acts and movements & Not promulgated to us as a text; known only in its participations & Immutable & Not a working juridical category; ground of all obligation & ST I-II 91.1; 93.1; 93.3\\
Natural law & God, legislating through the constitution of the rational creature; participation of the eternal law & Promulgated by being instilled in human minds, naturally knowable; formulated progressively by reason & First principles wholly immutable; secondary precepts stable \latin{ut in pluribus}; changeable by addition only & Not dispensable by any human authority; no custom against it obtains force of law & ST I-II 90.4 ad 1; 91.2; 94.2, 4--6; 97.3 ad 1; CIC cc.~24 §1, 199 1°, 1163 §2\\
Divine positive law & God, revealing in history (Old Law, New Law) & Given by divine acts and words; entrusted to the Church to guard and declare & Old Law's ceremonial and judicial precepts superseded within the one divine economy; New Law permanent & Not dispensable by the Church; the Church declares, does not constitute, its requirements & ST I-II 91.4--5; CIC cc.~748 §1, 750, 1075 §1, 1163 §2\\
Human law: ecclesiastical & The Church's competent legislative authority & Enacted and promulgated (cc.~7--8) or arising by legitimately observed custom approved by the legislator & Amendable, abrogable, subject to desuetude by contrary custom & Merely ecclesiastical law dispensable for a just cause within limits (cc.~85--93); custom \latin{contra legem} possible (cc.~24--26) & CIC cc.~11, 23--28, 85--93; ST I-II 95.1--2\\
Human law: civil & The civil community's legislative authority & Enacted, promulgated, or customary in the civil order & Amendable by its own order & Canon law defers to it in defined matters, \latin{quatenus iuri divino non sint contrariae} & CIC cc.~22, 98 §2, 1290; ST I-II 95.4; 96.1--4\\
\bottomrule
\end{longtable}
\endgroup

Three consequences of the map, stated now and tested against the Code in what follows. First, \emph{the boundary of dispensation and custom is the boundary of authorship}: what the Church did not author she cannot relax or watch lapse; what she authored she can. Second, \emph{the same written canon can carry more than one kind of law}: a canon may restate divine law in one paragraph and determine it ecclesiastically in the next (c.~1075 is the textbook case), so the classification attaches to the norm carried, not to the ink. Third, \emph{``positive'' is not a synonym for ``arbitrary''}: positive determinations bind in conscience when just, precisely because determination is one of the two modes of derivation from the natural law---though they bind, as q.~95 a.~2 says, with a force that is the enactment's own.
